\section{El recorrido, momento por momento}
El trámite avanza en cuatro momentos: 1. Reúne tus requisitos, 2. Propón sínodo y fecha, 3. Invitaciones al sínodo, 4. Examen y cierre. Es igual en todas las modalidades; las pantallas son de la modalidad \emph{Tesis}.
\subsection{Constancia de no adeudo y oficio de cotejo}
En el primer momento firmas la constancia de no adeudo de mobiliario y equipo y el oficio de cotejo.
\captura{tesis_coordinacion_01_firmas_no_adeudo}{«Posgrado: firmas» de la Coordinación: eliges los documentos, confirmas que los revisaste y firmas.}
\captura{tesis_coordinacion_02_firmas_cotejo}{La Coordinación firma el oficio de cotejo.}
\subsection{Invitaciones, acuse y notificación}
Con el examen programado firmas las invitaciones, el acuse y la notificación al director.
\captura{tesis_coordinacion_03_firmas_invitaciones}{La Coordinación firma las invitaciones, el acuse y la notificación; al firmar, cada integrante recibe la suya por correo.}
\subsection{Constancias e integración}
Al cierre firmas las constancias de revisores y de dirección y el formato de integración de expediente.
\captura{tesis_coordinacion_04_bandeja_cierre}{Constancias y formato de integración esperando la firma de la Coordinación.}
